\section{DBの概要}

\subsection{目的}

\begin{InyouJobunbox}
  \begin{itemize}
    \item 法第1条
  \end{itemize}
\end{InyouJobunbox}




\begin{fillblank}[年金2 \ 2021 \ 問題1(2)]

  \begin{lawstructure}
    \begin{lawarticle}{第一条}{目的}
      この法律は、（　　A　　）の進展、（　　B　　）の変化等の社会経済情勢の変化にかんがみ、
      事業主が従業員と給付の内容を約し、高齢期において従業員がその内容に基づいた給付を受けることが
      できるようにするため、確定給付企業年金について必要な事項を定め、
      国民の高齢期における所得の確保に係る（　　C　　）を支援し、
      もって（　　D　　）と相まって国民の生活の安定と
      （　　E　　）に寄与することを目的とする。
    \end{lawarticle}
  \end{lawstructure}

\end{fillblank}


\begin{answerbox}
 \begin{enumerate}[A:]
    \item 少子高齢化
    \item 産業構造
    \item 自主的な努力
    \item 公的年金の給付
    \item 福祉の向上
  \end{enumerate}
\end{answerbox}

\begin{fillblank}[数理人試験 \ 2023 \ 問題1 \ 設問1（抜粋）]

  \begin{lawstructure}
    \begin{lawarticle}{第一条}{目的}
      この法律は、少子高齢化の進展、（　　A　　）等の社会経済情勢の変化にかんがみ、
      事業主が従業員と（　　B　　）の内容を約し、高齢期において従業員がその内容に基づいた給付を受けること
      ができるようにするため、確定給付企業年金について必要な事項を定め、
      国民の高齢期における所得の確保に係る（　　C　　）を支援し、
      もって公的年金の給付と相まって（　　D　　）の安定と（　　E　　）の向上に寄与することを目的とする。
    \end{lawarticle}
  \end{lawstructure}

\end{fillblank}

\begin{answerbox}
 \begin{enumerate}[A:]
    \item 産業構造の変化
    \item 給付
    \item 自主的な努力
    \item 国民の生活
    \item 福祉
  \end{enumerate}
\end{answerbox}

\begin{fillblank}[年金1 \ H27 \ 問題1(2)]

  
  \noindent \textbf{◯確定給付企業年金法}
  \begin{lawstructure}
    \begin{lawarticle}{第一条}{目的}
      この法律は、少子高齢化の進展、（　　A　　）の変化等の社会経済情勢の変化にかんがみ、
      事業主が従業員と給付の内容を約し、高齢期において従業員がその内容に基づいた給付を受けること
      ができるようにするため、確定給付企業年金について必要な事項を定め、
      国民の高齢期における所得の確保に係る（　　B　　）を支援し、
      もって（　　C　　）と相まって国民の生活の安定と福祉の向上に寄与することを目的とする。
    \end{lawarticle}
  \end{lawstructure}

      
  \noindent \textbf{◯確定拠出年金法}
  \begin{lawstructure}
    \begin{lawarticle}{第一条}{目的}
      この法律は、少子高齢化の進展、（　　D　　）の多様化等の社会経済情勢の変化にかんがみ、
      個人又は事業主が拠出した資金を個人が（　　E　　）において運用の指図を行い、
      高齢期においてその結果に基づいた給付を受けることができるようにするため、
      確定拠出年金について必要な事項を定め、国民の高齢期における所得の確保に係る（　　B　　）を支援し、
      もって（　　C　　）と相まって国民の生活の安定と福祉の向上に寄与することを目的とする。
    \end{lawarticle}
  \end{lawstructure}

\end{fillblank}

\begin{answerbox}
 \begin{enumerate}[A:]
    \item 産業構造
    \item 自主的な努力
    \item 公的年金の給付
    \item 高齢期の生活
    \item 自己の責任
  \end{enumerate}
\end{answerbox}

\begin{explanationbox}
  DB法の第１条とDC法の第１条は非常に似ているのでセットで覚える。
\end{explanationbox}

\newpage

\subsection{各種用語の定義}

\begin{InyouJobunbox}
  \begin{itemize}
    \item 法第2条、第8条〜第10条、第94条
    \item 令第68条
  \end{itemize}
\end{InyouJobunbox}

\begin{fillblank}[数理人試験 \ 2023 \ 問題1 \ 設問1（抜粋）]

  \begin{lawstructure}
    \begin{lawarticle}{第二条}{定義}
      この法律において「確定給付企業年金」とは、厚生年金適用事業所の事業主が、単独で又は共同して、次章から第十三章までの規定に基づいて実施する年金制度をいう。
      \lawparagraph{2}{
        （略）
      }
      \lawparagraph{3}{
        この法律において「厚生年金保険の被保険者」とは、厚生年金保険の被保険者（厚生年金保険法第二条の五第一項（　　A　　）に規定する（　　A　　）厚生年金被保険者又は同項（　　B　　）に規定する（　　B　　）厚生年金被保険者に限る。）をいう。
      }
      \lawparagraph{4}{
        この法律において「企業年金基金」とは、前条の目的を達成するため、確定給付企業年金の加入者（以下「加入者」という。）に必要な給付を行うことを目的として、次章の規定に基づき設立された（　　C　　）をいう。
      }
    \end{lawarticle}
  \end{lawstructure}

\end{fillblank}

\begin{answerbox}
 \begin{enumerate}[A:]
    \item 第一号
    \item 第四号
    \item 社団
  \end{enumerate}
\end{answerbox}


\begin{fillblank}[自作 \ 2026/10/4]
  \begin{lawstructure}
    \noindent \textbf{◯DB法}
    \begin{lawarticle}{第九十四条}{福祉事業}
      基金は、第四章に規定する給付を行うほか、加入者等の（　　A　　）を増進するため、規約で定めるところにより、加入者等の（　　A　　）及び（　　B　　）に関する事業を行うことができる。
    \end{lawarticle}

    \noindent \textbf{◯DB令}
    \begin{lawarticle}{第六十八条}{会計の区分経理}
      加入者等の（　　A　　）及び（　　B　　）に関する事業を行う基金は、当該事業に係る会計を他の事業に係る会計と区分して経理しなければならない。
    \end{lawarticle}
  \end{lawstructure}
\end{fillblank}

\begin{answerbox}
 \begin{enumerate}[A:]
    \item 福祉
    \item 厚生
  \end{enumerate}
\end{answerbox}

\newpage

